\Article{17}
Die folgende Gleichung fasst somit die Bedingung zusammen, dass die zusammenwirkenden Kräfte die bewegliche Nadel nicht um die vertikale Achse drehen.
\[
0=-mT\sin u+fR^{-(n+1)}+f'R^{-(n+2)}+f''R^{-(n+3)}+\text{usw.}-\theta(u-N)
\]

Der Wert von $N$ lässt sich leicht so einstellen, dass er, wenn er nicht genau $=0$ ist, wenigstens sehr klein bleibt. Auch $u$ bleibt bei den hier betrachteten Versuchen innerhalb enger \SourceLeaf{109}Grenzen. Daher kann man den Term $\theta(u-N)$ ohne Furcht vor einem merklichen Fehler durch $\theta\sin(u-N)$ ersetzen, zumal $\frac{\theta}{mT}$ ein sehr kleiner Bruch ist. Sei $u^0$ der Wert von $u$, der dem Gleichgewicht der ersten Nadel in Abwesenheit der zweiten entspricht. Es gelte also
\[
mT\sin u^0+\theta\sin(u^0-N)=0
\]
Daraus folgt leicht
\[
mT\sin u+\theta\sin(u-N)=(mT\cos u^0+\theta\cos(u^0-N))\sin(u-u^0)
\]
Hier können wir den ersten Faktor unbedenklich durch $mT+\theta$ ersetzen. Unsere Gleichung nimmt damit die Form an
\[
(mT+\theta)\sin(u-u^0)=fR^{-(n+1)}+f'R^{-(n+2)}+f''R^{-(n+3)}+\text{usw.}
\]
Behalten wir hier nur den ersten Term $fR^{-(n+1)}$ bei, so ergibt sich die Lösung unmittelbar. Wir erhalten
\[
\begin{aligned}
&\operatorname{tang}(u-u^0)\\
&=\frac{\begin{gathered}mM(n\cos(\psi-U)\sin(\psi-u^0)\\+\sin(\psi-U)\cos(\psi-u^0))R^{-(n+1)}\end{gathered}}{\begin{gathered}mT+\theta+mM(n\cos(\psi-U)\cos(\psi-u^0)\\-\sin(\psi-U)\sin(\psi-u^0))R^{-(n+1)}\end{gathered}}
\end{aligned}
\]
Mit demselben Recht können wir im Nenner den Teil weglassen, der den Faktor $R^{-(n+1)}$ enthält, oder setzen
\[
\begin{aligned}
\operatorname{tang}(u-u^0)&=\frac{mM}{mT+\theta}(n\cos(\psi-U)\sin(\psi-u^0)\\
&\qquad+\sin(\psi-U)\cos(\psi-u^0))R^{-(n+1)}\\
&=FR^{-(n+1)}
\end{aligned}
\]

Wollen wir dagegen die weiteren Terme berücksichtigen, so ist klar, dass sich $\operatorname{tang}(u-u^0)$ in eine Reihe folgender Form entwickeln lässt
\[
\operatorname{tang}(u-u^0)=FR^{-(n+1)}+F'R^{-(n+2)}+F''R^{-(n+3)}+\text{usw.}
\]
Eine kurze Überlegung zeigt, dass die Koeffizienten $F,F',F''$ usw. bis einschließlich zum Koeffizienten der Potenz $R^{-(2n+1)}$ der Reihe nach aus
\[
\frac{f}{mT+\theta},\quad\frac{f'}{mT+\theta},\quad\frac{f''}{mT+\theta}\quad\text{usw.}
\]
entstehen, indem man $u$ durch $u^0$ ersetzt. Vom nächsten Term an kommen neue Beiträge hinzu, die wir für unseren Zweck jedoch nicht genauer verfolgen müssen. Offenbar lässt sich übrigens $u-u^0$ in eine Reihe ähnlicher Form entwickeln, die sogar bis zur Potenz $R^{-(3n+2}$ mit der Reihe für $\operatorname{tang}(u-u^0)$ übereinstimmt.

